\section{Captive：360度全景生成}

\subsection{问题与挑战}

360度全景图像的生成面临一个严峻的数据问题：真实的360度全景数据\textbf{非常有限}。如何在有限数据下训练高质量的全景生成模型？

\subsection{核心思想：立方体贴图表示}

Captive 的关键洞察是将360度全景表示为\textbf{立方体贴图}（cube map），即6个面的透视图像。每个单独的面都是标准透视图像，属于互联网图像的正常分布范围。

\begin{importantbox}{Captive 的架构设计}
\textbf{编码}：将6张立方体面沿通道维度堆叠，输入2D编码器。

\textbf{生成}：几乎完全复用预训练 Stable Diffusion 模型的架构。由于预训练模型已具备生成任意透视图像的能力，\textbf{冻结现有层}，仅训练新增的注意力层。

\textbf{新增层的作用}：学习6张图像之间的空间关系，即立方体各面如何在3D空间中关联。

\textbf{数据效率}：由于核心生成能力来自预训练的大规模2D模型，仅需\textbf{数万张}全景图像即可训练出高质量的全景生成模型。
\end{importantbox}

\subsection{分布外泛化}

尽管训练数据仅限于自然场景的全景图，Captive 能够生成超出训练分布的内容。这得益于冻结了2D基础模型——它保留了生成任意类型透视图像的能力，而新增的注意力层只需学习视角间的空间关系。

\begin{knowledgebox}{冻结基础模型 + 轻量适配层的范式}
Captive 的成功体现了一个重要范式：冻结强大的预训练基础模型，只训练少量的适配层来学习新的约束（如多视角一致性、360度空间关系）。这样既能利用基础模型在大规模数据上学到的强大生成先验，又能在小数据集上快速适配新任务，且保持了基础模型的分布外泛化能力。
\end{knowledgebox}

\subsection{本章小结}

Captive 通过立方体贴图表示和冻结基础模型的策略，在有限全景数据下实现了高质量的360度全景生成，并展现了良好的分布外泛化能力。

